\documentclass[12pt,a4paper,UTF8]{ctexbook}
\usepackage{geometry,amsmath,amssymb,amsthm,booktabs,hyperref}
\geometry{left=3.0cm,right=2.5cm,top=2.8cm,bottom=2.6cm}
\hypersetup{colorlinks=true,linkcolor=blue,citecolor=blue,urlcolor=blue}
\newtheorem{definition}{定义}[chapter]
\newtheorem{proposition}{命题}[chapter]
\newtheorem{theorem}{定理}[chapter]
\title{面向金融问答的时点约束义务感知检索规划研究}
\author{匿名作者}
\date{2026年9月}
\begin{document}
\begin{titlepage}
\centering
\vspace*{1.5cm}{\zihao{2}\bfseries 上海财经大学本科毕业论文（初稿）\\[1.5cm]}
{\zihao{1}\bfseries 面向金融问答的时点约束义务感知检索规划研究\\[2cm]}
\begin{tabular}{rl}学院：&金融学院\\[0.4cm]专业：&金融学（数据与智能方向）\\[0.4cm]作者：&匿名\\[0.4cm]指导教师：&待定\end{tabular}
\vfill 2026年9月
\end{titlepage}
\frontmatter
\chapter*{摘要}
金融问答并不只是从语料库中找出最相关的一句话。一个问题可能同时要求多个实体、财务年度、季度以及申报类型，并且只有在问题时点之前公开的文件才具备证据资格。本文提出 FinPlan-RAG，一种非预言式的时点约束义务感知检索规划方法。系统将问题解析为 $(实体,财务年度,申报类型,财务期间)$ 四元组义务，针对尚未解决的义务逐项检索，在固定查询预算下仅接纳时点有效的文档，并在证据闭合前保持保守状态。

本文首先说明研究缺口：单一相关性分数不能保证停止决策所需的证据集合已经闭合。随后给出两个世界具有相同可见分数但需要相反动作的构造性证明，并证明去掉时点过滤后可以发生未来信息泄漏。在本地可复现实验中，我们使用 20 个随机种子、7 类场景、每类 40 个样本和统一预算 10 进行 5,600 个 case 的机制检验。FinPlan-RAG 的闭合率为 0.366，而 single-shot 基线为 0.101；平均查询次数为 8.87，而基线为 10。去掉反证模块后闭合率上升至 0.511，但反证召回率由 0.605 降至 0.315，过度断言率由 0.013 升至 0.035；去掉时点过滤后未来文档泄漏率达到 0.502。上述结果仅支持受控机制结论，不构成真实财报问答或顶刊 SOTA 结论。

本文同时提供面向 Codex、Claude Code 等代码 Agent 的统一 JSON 工具协议和标准输入输出适配器。研究结果、协议哈希和测试命令均在代码仓库中公开，以便后续使用冻结的金融文件数据、共享阅读器和强基线进行独立复核。

\noindent\textbf{关键词：} 金融问答；检索增强生成；时点约束；义务状态；证据闭合；反证；代码 Agent
\tableofcontents
\mainmatter
\chapter{绪论}
\section{研究背景}
金融披露具有明确的实体、年度、季度、申报类型和公开时间。传统 RAG 通常将检索看作排序问题，返回若干高相关段落，却没有显式表示“还缺少哪一份申报文件”。这种隐含状态会造成遗漏第二个实体或期间、把后续修订文件带入历史问题、以及忽视限制性披露等风险。
\section{研究目标与假设}
本文检验：在信息、阅读器和查询预算相同的条件下，显式义务状态是否能够改善下一步检索和停止决策。核心假设 H1 为：义务感知状态提高证据闭合率并减少无效查询；显式反证和时点过滤牺牲部分覆盖率，但降低过度断言和未来信息泄漏。
\section{研究缺口与贡献}
现有相关性排序、稀疏/稠密检索和长上下文记忆基准解决了“哪些文本相似”问题，但没有保证“所有被问题要求的文件都已被覆盖”。本文形式化金融检索义务，给出相关性分数不足以支撑停止决策的证明，实现非预言式规划器，开展基线、消融及时点压力测试，并提供代码 Agent 统一接口。
本文沿用 SQCAD 的生命周期分层思想，将候选提出、证据资格判断和阅读器生成保持分离，并将其具体化为金融申报义务与时点约束。
\chapter{理论基础与研究缺口证明}
\section{义务与闭合}
对问题 $q$，令 $O(q)$ 为义务集合，义务元素记为 $o=(e,y,f,p)$；令 $D_{\le c}$ 表示公开时间不晚于截止时点 $c$ 的文件。
\begin{definition}[证据闭合]
当每个具有已知文档编号的义务都至少被一个时点有效文档覆盖，且所需字段的投票结果达到解析阈值时，称证据状态闭合。闭合是检索过程性质，不等于答案正确。
\end{definition}
\begin{proposition}[同分数反例]
若存在两个潜在世界具有相同可见历史和分数，但对未解决义务的最优动作相反，则任意仅依赖该分数的随机规则都具有非零最坏世界后悔。设正、负世界的动作边际分别为 $m_+$ 和 $m_-$，规则以概率 $p$ 选择正动作，则
\[
\max\{(1-p)m_+,pm_-\}\geq\frac{m_+m_-}{m_++m_-}。
\]
\end{proposition}
证明：正世界选择负动作的概率为 $1-p$，负世界选择正动作的概率为 $p$；平衡两项损失下界即得结论。
\begin{theorem}[义务状态价值分解]
令 $x_t=(U_t,E_t,c)$ 包含未解决义务集合、证据状态和截止时间，$a_t$ 为义务定向检索或停止动作。若即时效用为 $g(x_t,a_t)$，转移与观测核为 $K_a(dx_{t+1},do_{t+1}\mid x_t)$，则有限时域动作价值满足
\[
Q_t(x,a)=\mathbb{E}[g(x,a)]+\gamma\mathbb{E}_{K_a}V_{t+1}(x'),\qquad V_t(x)=\max_aQ_t(x,a)。
\]
两动作 $a,b$ 的生命周期差异为即时差异与延续差异之和。如果分数摘要 $S(x)$ 能支持最优分数策略，则 $Q_t(x,a)-Q_t(x,b)$ 必须关于 $S(x)$ 可测；否则存在同分数但差异符号相反的世界。
\end{theorem}
证明：对下一状态和观测进行条件化得到 Bellman 递推；相减得到差异分解。若差异不关于 $S$ 可测，则存在同分数且符号相反的状态，代入前述命题得到严格正的最坏世界后悔。
\section{时点安全性}
\begin{theorem}[未来信息过滤]
若规划器只将 $available\_at\leq c$ 的片段写入证据状态，则输出上下文不包含未来文件。若去掉该谓词，则可以构造未来文件相关性更高而被最大分数选择消费的 case。
\end{theorem}
\chapter{FinPlan-RAG 算法设计}
\section{总体框架}
系统流程为：问题解析 $\rightarrow$ 义务生成 $\rightarrow$ 时点过滤 $\rightarrow$ 义务定向 BM25 检索 $\rightarrow$ 证据状态聚合 $\rightarrow$ 闭合检查 $\rightarrow$ 阅读器回答或拒答。数据结构包括 DocumentObligation、Chunk、RetrievalResult 和 EvidenceState。
\section{期间解析与义务生成}
系统将 Q1--Q3 映射为 10-Q，将 Q4、FY 和 annual 映射为 10-K/FY；“上半年”映射为 Q2，“下半年”映射为 Q3 与 Q4。解析结果与可见申报元数据求交，输出稳定排序的义务元组。
\section{检索与停止}
对每个未解决义务 $o$，构造 $q\oplus o$ 查询，并要求片段包含实体和年份词。BM25 只从截止时点之前的片段中选择一个最高分结果。状态写为
\[
E_t=(O(q),C_t,n_t,closed_t,F_t)，
\]
其中 $C_t$ 为已覆盖文档集合，$n_t$ 为查询次数，$F_t$ 为未来文档诊断集合。闭合前阅读器可以拒答或请求继续检索，不能把闭合当作答案正确性的替代品。
\section{代码 Agent 接口}
仓库新增 \texttt{finplan-rag-agent --schema} 命令，输出 OpenAI-compatible function schema；默认模式按 JSONL 读取请求并输出 JSONL 响应。响应包含义务、覆盖文档、查询次数、闭合标记、未来文档诊断、选中的时点有效证据片段以及 \texttt{answer\_with\_reader} 或 \texttt{abstain\_or\_retrieve\_more} 建议。
\chapter{实验设计与结果}
\section{实验协议}
受控实验使用 20 个种子、7 个场景、每场景每种子 40 个 case，总计 5,600 个 case，预算为 10。所有方法共享路径、文件、截止时点、解析器和预算。效用函数预先固定为：正确 $=1$、拒答 $=-0.15$、其他错误 $=-1$、过度断言额外 $=-1$、每次查询 $=-0.03$。完整结果文件的 SHA-256 为 \texttt{E5D207A2...42CAA28}。
\section{主结果}
\begin{table}[ht]\centering\caption{受控机制实验均值。除过度断言、查询次数和未来泄漏外，数值越大越好。}\begin{tabular}{lrrrrrrr}\toprule
方法 & 准确率 & 闭合率 & 过度断言 & 反证召回 & 查询数 & 效用 & 未来泄漏\\\midrule
Single-shot & .066 & .101 & .004 & .651 & 10.00 & -.407 & .000\\
固定分解 & .175 & .293 & .014 & .658 & 10.00 & -.363 & .000\\
通用自适应 & .193 & .343 & .015 & .632 & 10.00 & -.371 & .000\\
FinPlan 状态 & .191 & .366 & .013 & .605 & 8.87 & -.358 & .000\\\bottomrule\end{tabular}\end{table}
相对 single-shot，FinPlan 的闭合率提高 0.2645，查询次数减少 1.1279。与通用自适应相比，准确率差异为 -0.0020，说明当前证据主要支持流程和停止改进，而非普遍答案准确率优势。
\section{消融与压力测试}
去掉反证后闭合率为 0.511，查询次数为 7.92，但反证召回由 0.605 降为 0.315，过度断言由 0.013 升为 0.035。去掉自适应停止后始终使用 10 次查询；去掉时点过滤后未来文档泄漏率为 0.502，效用降至 -0.438。
\begin{table}[ht]\centering\small\caption{同一 5,600-case 合同下的组件消融。}\begin{tabular}{lrrrrrr}\toprule
变体 & 闭合率 & 过度断言 & 反证召回 & 查询数 & 效用 & 未来泄漏\\\midrule
FinPlan 状态 & .366 & .013 & .605 & 8.87 & -.358 & .000\\
去掉反证 & .511 & .035 & .315 & 7.92 & -.284 & .000\\
去掉自适应停止 & .337 & .012 & .606 & 10.00 & -.370 & .000\\
去掉时点过滤 & .433 & .021 & .578 & 8.57 & -.438 & .502\\\bottomrule\end{tabular}\end{table}
\section{预算敏感性}
在相同 20 个种子和 7 个场景下，将预算设置为 1、2、4、6、10、16。FinPlan 闭合率依次为 0.000、0.000、0.029、0.182、0.366 和 0.455，反证召回率由 0.315 升至 0.633；由于查询成本持续累积，效用由 -0.180 降至 -0.451。因此预算 10 只是可复现的工作点，不应被解释为合同之外的全局最优。完整结果与每次运行的哈希见 \texttt{paper/results/budget\_sensitivity.json}。
\section{真实数据外部效度审计}
为避免把受控世界误写成金融证据，本文单独审计了早期 FinPlanRAG 实验的只读冻结结果。公开 LOFin 全集包含 1,572 道题，预算为 4。FinPlan v9 的文件级闭合率为 0.8919，通用期间自适应基线为 0.8925，HiREC 风格基线为 0.8588，single-shot 为 0.6425；对应错误文档率分别为 0.3108、0.4742、0.6244 和 0.7983。扩展配对检验支持其相对通用、HiREC 风格和 single-shot 的错误文档率优势，并支持多文档分层的闭合优势，但与强期间元数据基线的全集闭合并无优势。更关键的是，义务精确匹配仅为 0.4637，且 metadata-fill 对照与级联策略完全持平。因此该结果支持“义务表示层”的解释，不支持独立的动作策略优势或答案准确率优势。

中国 FinGLM 审计覆盖 1,829 道题和 11,587 份年报。FinPlan v9 与 metadata-fill 的文件闭合率均为 0.9978、错误文档率均为 0.0126；期间元数据基线为 0.9858/0.0120，通用自适应为 0.9885/0.2712，single-shot 为 0.7594/0.8087。该轨道使用根据题目元数据推导的单文档金标，只评估文件闭合，不评估原始答案，因此仅能说明多年义务解析收益。SEC 点时 pilot 含 17 条探索 lineage 和 12 条留出 lineage；在留出审计中，去掉时间过滤后未来泄漏为 5/12，而保留时间过滤为 0/12，配对检验为描述性结果（$p=0.0625$）。中国 mixed-period pilot 仅有 5 道题（冻结 3 道）：期间感知规划冻结闭合 3/3，实体—年份规划闭合 0/3；其阅读器是正则表达式字段审计器，不构成真实 LLM 答案评测，也不作显著性主张。四条轨道的哈希和外部只读路径见 \texttt{paper/results/real\_data\_summary.json} 与 \texttt{paper/real\_data\_provenance.md}。
\chapter{讨论与局限}
实验只支持受控机制层结论。合成世界不能替代真实 SEC/交易所文件，闭合率不能替代答案正确率，BM25 也不能代表强稠密检索器。当前实现没有声称训练新模型、没有把 SQCAD 的结果当作 FinPlan 的结果，也没有声称顶刊 SOTA。下一阶段应冻结真实财报数据、共享阅读器、时间切分和强基线，并报告答案正确率、证据闭合、未来泄漏、过度断言和成本的联合 Pareto 前沿。
\chapter{结论}
本文将金融 RAG 重新表述为截止时点下的义务完成问题。理论反例表明，仅凭相关性分数无法保证停止决策；实现与受控实验显示，显式状态提高闭合并减少查询，但反证和时点约束会引入可测成本。代码 Agent 适配器使 Codex、Claude Code 等工具能够在同一可审计契约下调用系统。
\backmatter
\chapter*{参考文献}
\begin{thebibliography}{9}
\bibitem{lewis2020rag} Lewis, P. 等. Retrieval-augmented generation for knowledge-intensive NLP tasks. NeurIPS, 2020.
\bibitem{maharana2024locomo} Maharana, A. 等. LoCoMo: Long-context conversation memory benchmark. ACL, 2024.
\bibitem{wu2024longmemeval} Wu, L. 等. LongMemEval: Benchmarking long-term memory of large language models. arXiv, 2024.
\end{thebibliography}
\chapter*{附录：证据与复现清单}
表中数字来自 \texttt{paper/results/controlled\_summary.json}，源文件哈希记录在同目录 README。运行命令为 \texttt{PYTHONPATH=src python -m finplan\_rag.controlled\_planning --seeds 20 --cases-per-world 40 --budget 10}；测试命令为 \texttt{python -m pytest -q}。大规模语料、模型权重和真实数据遵循仓库 \texttt{DATA\_STORAGE.md}，不写入 Git。
\end{document}
